\section{Introduction}\label{sec:introduction}

A function \(f:\R\to\R\) is \(\PF_r\) when every translation-kernel minor
\[
 \det[f(x_i-y_j)]_{i,j=1}^k
\]
is nonnegative for \(1\leq k\leq r\) and increasing node tuples. We call the
property strict when the determinant is positive whenever both tuples are
strictly increasing. This paper determines the exact finite P\'olya-frequency
order of the Riemann kernel; see \cite{belton} for modern background and
Schoenberg's characterization. Related finite-order and Riemann-kernel work is
placed in \secref{sec:related-work}.

\begin{theorem}[Exact global PF order]\label{thm:main}
The exact global P\'olya-frequency order of \(\Phi\) is four. More precisely:
\begin{enumerate}
\item for every \(1\leq k\leq4\) and strictly increasing real tuples
\(x_1<\cdots<x_k\) and \(y_1<\cdots<y_k\),
\[
 \det[\Phi(x_i-y_j)]_{i,j=1}^k>0.
\]
\item with \(h=211/2000\),
\[
 \det[\Phi((i-j)h)]_{i,j=0}^4<0.
\]
\end{enumerate}
\end{theorem}

The two theorem inputs are deliberately separated. The positive input is the
universal strict-order theorem through four. The negative input is the direct
rational order-five witness in part~(2), including proved strict ordering of
its nodes and the signed Toeplitz-to-translation-matrix identification.

The proof has two one-variable inputs and one exact transport
mechanism. First, global inequalities \(q>0\) and \(F_2>0\) imply strict
\(\PF_3\) and positivity of a three-point functional \(\Lambda\). Second, the
fully confluent order-four invariant \(\Cfour\) is globally positive. Exact
quotient identities reduce strict \(\PF_4\) to the monotonicity of a scalar
three-point function \(\Psi\). On the actual image of an increasing curvature
coordinate, the derivative of \(\Psi\) becomes a positive integral of
\(\Cfour\) against a deterministic closed cumulative gap.

The finite part of the positive input consists of exact rational Bernstein
coefficients on fixed algebraic domains, followed by analytic bounds for all
remaining theta modes; \appref{app:tail} lists the transformed target
polynomials, domains, degrees, minimum coefficients, derivative envelopes,
and the surviving relative margins. The transport algebra is proved in the
text, and the order-five
calculation uses rational enclosures and elementary series bounds.
